% RQ1.4b / RQ1.5b multi-node placement — server-rack stylisation.
% Each Kubernetes node is drawn as a stylised server chassis (rounded
% chassis box, drive-bay slats, power LED) so the reader sees six distinct
% physical machines. The pod is nested inside its chassis. gRPC arrows
% cross the air-gap between adjacent chassis.
\begin{tikzpicture}[
    font=\small,
    cluster/.style={
        draw=black!50,
        fill=white,
        rounded corners=10pt,
        thick,
        inner xsep=18pt,
        inner ysep=22pt
    },
    chassis/.style={
        draw=black!70,
        fill=black!8,
        line width=1.2pt,
        rounded corners=3pt,
        minimum width=2.6cm,
        minimum height=3.1cm
    },
    chassisclient/.style={
        draw=orange!75!black,
        fill=orange!10,
        line width=1.2pt,
        rounded corners=3pt,
        minimum width=2.6cm,
        minimum height=3.1cm
    },
    pod/.style={
        draw=blue!65!black,
        fill=blue!18,
        rounded corners=4pt,
        minimum width=2.05cm,
        minimum height=1.0cm,
        align=center,
        thick
    },
    clientpod/.style={
        draw=orange!75!black,
        fill=orange!22,
        rounded corners=4pt,
        minimum width=2.05cm,
        minimum height=1.0cm,
        align=center,
        thick
    },
    svc/.style={
        draw=green!45!black,
        fill=green!15,
        rounded corners=3pt,
        minimum width=1.7cm,
        minimum height=0.45cm,
        align=center,
        thick,
        font=\tiny
    },
    arrow/.style={
        -{Latex[length=2.6mm]},
        very thick,
        draw=red!75!black
    },
    returnarrow/.style={
        -{Latex[length=2.4mm]},
        thick,
        draw=black!55,
        dashed
    },
    smalllabel/.style={
        font=\scriptsize,
        text=black!65,
        align=center
    },
    transferlabel/.style={
        font=\scriptsize\bfseries,
        text=red!75!black,
        align=center
    },
    nodelabel/.style={
        font=\small\bfseries,
        text=black!75
    },
    skulabel/.style={
        font=\tiny\itshape,
        text=black!55
    }
]

% --- Place 6 server chassis ---
\node[chassisclient] (n0) at (0, 0) {};
\node[chassis,       right=1.3cm of n0] (n1) {};
\node[chassis,       right=1.3cm of n1] (n2) {};
\node[chassis,       right=1.3cm of n2] (n3) {};
\node[chassis,       right=1.3cm of n3] (n4) {};
\node[chassis,       right=1.3cm of n4] (n5) {};

% --- Helper: draw drive-bay slats + LED + SKU label inside each chassis ---
\foreach \n in {n0, n1, n2, n3, n4, n5} {
    % Drive-bay slats (3 horizontal lines near the top)
    \foreach \dy in {0.18, 0.32, 0.46} {
        \draw[draw=black!45, line width=0.55pt]
            ([xshift=4pt, yshift=-\dy cm]\n.north west) --
            ([xshift=-4pt, yshift=-\dy cm]\n.north east);
    }
    % Power LED (small green dot, top-right)
    \fill[green!55!black]
        ([xshift=-5pt, yshift=-7pt]\n.north east) circle (1.2pt);
    \draw[draw=black!45, line width=0.4pt]
        ([xshift=-5pt, yshift=-7pt]\n.north east) circle (1.2pt);
    % SKU label centred between slats and content
    \node[skulabel] at ([yshift=-0.66cm]\n.north) {Standard\_D8s\_v3};
}

% --- Pods + services nested inside each chassis ---
\node[clientpod] at ([yshift=-0.55cm]n0.center) (client) {Benchmark\\client pod};
\node[pod]       at ([yshift=-0.30cm]n1.center) (p1)     {Pod 1\\\texttt{stem+layer1}};
\node[pod]       at ([yshift=-0.30cm]n2.center) (p2)     {Pod 2\\\texttt{layer2}};
\node[pod]       at ([yshift=-0.30cm]n3.center) (p3)     {Pod 3\\\texttt{layer3}};
\node[pod]       at ([yshift=-0.30cm]n4.center) (p4)     {Pod 4\\\texttt{layer4}};
\node[pod]       at ([yshift=-0.30cm]n5.center) (p5)     {Pod 5\\\texttt{avgpool+fc}};

\node[svc, below=0.18cm of p1] (s1) {Service 1};
\node[svc, below=0.18cm of p2] (s2) {Service 2};
\node[svc, below=0.18cm of p3] (s3) {Service 3};
\node[svc, below=0.18cm of p4] (s4) {Service 4};
\node[svc, below=0.18cm of p5] (s5) {Service 5};

% --- Node labels above each chassis ---
\node[nodelabel, above=2pt of n0.north] {Node 0 (client)};
\node[nodelabel, above=2pt of n1.north] {Node 1};
\node[nodelabel, above=2pt of n2.north] {Node 2};
\node[nodelabel, above=2pt of n3.north] {Node 3};
\node[nodelabel, above=2pt of n4.north] {Node 4};
\node[nodelabel, above=2pt of n5.north] {Node 5};

% --- Cluster background ---
\begin{pgfonlayer}{background}
\node[
    cluster,
    fit=(n0)(n1)(n2)(n3)(n4)(n5)(s1)(s2)(s3)(s4)(s5),
    label={[font=\small\bfseries, text=black!75]above:%
        AKS cluster \textemdash{} 6 \texttt{Standard\_D8s\_v3} VMs, single VNet, anti-affinity enforced}
] (clusterbox) {};
\end{pgfonlayer}

% --- Inter-node gRPC arrows that visibly cross the air-gap between chassis ---
\draw[arrow] (n0.east) -- node[transferlabel, above] {gRPC} (n1.west);
\draw[arrow] (n1.east) -- node[transferlabel, above] {784 KiB} (n2.west);
\draw[arrow] (n2.east) -- node[transferlabel, above] {392 KiB} (n3.west);
\draw[arrow] (n3.east) -- node[transferlabel, above] {196 KiB} (n4.west);
\draw[arrow] (n4.east) -- node[transferlabel, above] {98 KiB}  (n5.west);

% Pod-to-service association lines
\draw[black!35, thick] (s1) -- (p1);
\draw[black!35, thick] (s2) -- (p2);
\draw[black!35, thick] (s3) -- (p3);
\draw[black!35, thick] (s4) -- (p4);
\draw[black!35, thick] (s5) -- (p5);

% Return path (curves below all chassis)
\draw[returnarrow]
    (n5.south) .. controls +(0.4,-1.6) and +(0,-1.8) ..
    node[smalllabel, below, pos=0.85] {final output returned to caller}
    (n0.south);

\end{tikzpicture}%
